\chapter{Hướng dẫn chạy và phân công chi tiết}
\label{chap:phu-luc}

\section{Chạy ứng dụng trên môi trường phát triển}
Cần Node.js 24, PostgreSQL và thông tin kết nối riêng của môi trường.
Các bước dưới đây dùng PowerShell; thực hiện từ thư mục gốc repository.
Không ghi đè file env đang sử dụng. Nếu file đã có, kiểm tra và cập nhật
các biến cần thiết theo env mẫu.

\begin{lstlisting}[caption={Chuẩn bị backend và dữ liệu menu}]
Set-Location apps/api
npm ci
# Chi sao chep khi chua co file .env
Copy-Item .env.example .env
# Dien DATABASE_URL va JWT_SECRET rieng trong .env.
# Tao database PostgreSQL truoc khi chay migrate.
npm run prisma:generate
npm run prisma:deploy
npm run db:seed
npm run start:dev
\end{lstlisting}
Mở terminal khác để chạy frontend:
\begin{lstlisting}[caption={Chạy frontend từ gốc repository}]
Set-Location apps/web
npm ci
npm run dev
\end{lstlisting}
Frontend mặc định tại \url{http://localhost:3000}, API có tiền tố
\url{http://localhost:3001/api/v1}. Env frontend theo
\sourcefile{.env.example}: cần kiểm tra \code{NEXT_PUBLIC_API_URL} khớp
địa chỉ API. Kiểm tra \code{/health}, menu và đăng ký một tài khoản
CUSTOMER trước khi demo. Không dùng seed E2E làm quy trình bàn giao
chính thức hoặc coi seed catalog là đã tạo sẵn tài khoản Barista.

Compose có thể khởi chạy bằng \code{docker compose up --build} sau khi
chuẩn bị env gốc. Service migrate hiện không seed catalog: phải thực
hiện bước seed với cùng database qua môi trường API đã cấu hình, rồi
kiểm tra menu. Hướng dẫn Compose sạch và tài khoản demo là hạng mục
cần được nhóm kiểm chứng khi bàn giao.

\section{Lệnh kiểm tra ứng dụng}
\begin{lstlisting}[caption={Các lệnh kiểm tra dùng trong từng ứng dụng}]
# Trong apps/api
npm run lint
npm run typecheck
npm test
npm run test:e2e
# DB integration can SMOKE_DATABASE_URL cua DB test rieng.
npm run test:integration
npm run build

# Trong apps/web
npm run lint
npm run typecheck
npm test
npm run test:e2e
npm run build
\end{lstlisting}
DB integration và browser E2E phải theo cấu hình fixture của repo và
hướng dẫn \sourcefile{apps/api/README.md},
\sourcefile{tests/e2e/README.md}. Không trỏ các bài ghi dữ liệu vào DB
đang dùng của ứng dụng. Kết quả của các lệnh này cần được lưu cùng mốc
source khi thực sự chạy, không mặc định lấy từ bảng số liệu lịch sử.

\section{Danh mục thẻ phân công và reviewer}
\label{sec:task-assignment}
Bảng trích từ bản đối chiếu phân công ngày 03/11/2026~\cite{assignment}.
Tên ngắn tương ứng các thành viên đầy đủ ở Bảng~\ref{tab:assignment}.
Cột reviewer thể hiện người được giao kiểm tra, chưa xác nhận đã approve.
Tên đầu việc giữ theo nội dung phân công; kỹ thuật cũ trên thẻ cần đọc
cùng thiết kế/hiện trạng mới trong báo cáo.

{\footnotesize\setstretch{1.08}
\begin{longtable}{@{}L{0.15\textwidth}L{0.50\textwidth}L{0.12\textwidth}L{0.13\textwidth}@{}}
\caption{Danh mục 40 thẻ, người phụ trách và người kiểm tra}\label{tab:task-assignment}\\
\toprule\textbf{Mã thẻ} & \textbf{Đầu việc được giao} & \textbf{Phụ trách} & \textbf{Reviewer} \\\midrule\endfirsthead
\toprule\textbf{Mã thẻ} & \textbf{Đầu việc được giao} & \textbf{Phụ trách} & \textbf{Reviewer} \\\midrule\endhead
\bottomrule\endfoot
A1-1.1 & Repository và khung API NestJS & Khang & Huy \\
A1-1.2 & Khởi tạo Next.js và kết nối backend & Khang & Huy \\
A1-1.4 & Env mẫu, lint/build và hướng dẫn môi trường & Khang & Huy \\
A1-2.1 & Contract Product và GET /products & Khang & Kiệt \\
A1-10.1 & Thiết kế Order State Machine & Khang & Kiệt \\
SCRUM-S1 & Review/Retro và hồ sơ Sprint & Khang & Khải \\
A1-6.1 & POST /orders, giá server và PENDING & Khang & Hoàng \\
A1-6.2 & Order snapshot, transaction và rollback & Khang & Hoàng \\
A1-9.1 & API lịch sử và chi tiết đơn của tôi & Khang & Huy \\
SCRUM-S2 & Review/Retro và hồ sơ Sprint & Khang & Khải \\
A1-9.5 & Docker Compose và README bàn giao & Khang & Hoàng \\
SCRUM-S3 & Review/Retro và hồ sơ Sprint & Khang & Khải \\
A1-1.3 & Kết nối PostgreSQL/Prisma và dữ liệu mẫu & Hoàng & Khang \\
A1-2.2 & Seed menu và API GET /products & Hoàng & Khải \\
A1-2.3 & Kiểm thử API danh sách sản phẩm & Hoàng & Khải \\
A1-3.1 & Lưới menu và thẻ đồ uống mobile-first & Hoàng & Khải \\
A1-3.2 & Tích hợp API thật vào Menu & Hoàng & Khải \\
A1-3.3 & Loading, empty, error và retry của Menu & Hoàng & Khải \\
A1-4.1 & Product Detail API và tùy chọn hợp lệ & Hoàng & Khang \\
A1-4.3 & Kiểm thử API tùy chọn và giá sản phẩm & Hoàng & Khang \\
A1-10.4 & Inventory reserve/consume/release và concurrency & Hoàng & Khang \\
A1-9.3 & Migration/seed và dữ liệu demo sạch & Hoàng & Khải \\
A1-7.1 & Register/Login, bcrypt và JWT & Huy & Khang \\
A1-7.2 & Auth UI, phiên, logout và 401 & Huy & Khải \\
A1-6.3 & JWT, validation và chủ sở hữu đơn & Huy & Khải \\
A1-7.3 & Bảo vệ API đặt đơn và token an toàn & Huy & Khải \\
A1-10.6 & Loyalty ledger/balance sau PAID & Huy & Khang \\
A1-9.4 & Tài khoản demo, bảo mật và hướng dẫn test & Huy & Khang \\
A1-4.2 & Detail UI, size/topping và giá tạm tính & Khải & Hoàng \\
A1-5.1 & Cart store thêm/sửa/xóa và persist & Khải & Hoàng \\
A1-5.2 & Định danh dòng giỏ và tính tổng & Khải & Hoàng \\
A1-5.3 & Kiểm thử phép tính và khôi phục giỏ & Khải & Hoàng \\
A1-8.2 & Checkout UI và xử lý payment outcome & Khải & Khang \\
A1-9.2 & Confirmation và History/Detail UI & Khải & Hoàng \\
A1-10.5 & Coupon validation và retry an toàn & Kiệt & Huy \\
A1-8.1 & Mock gateway và POST /payments & Kiệt & Khang \\
A1-8.3 & PAID/PAYMENT\_FAILED qua OrdersService & Kiệt & Khang \\
A1-10.3 & Payment idempotency và chống double-charge & Kiệt & Khang \\
A1-10.2 & Hiện thực full Order State Machine & Kiệt & Khang \\
EXT-UC10 & Staff UI queue và cập nhật trạng thái đơn & Kiệt & Huy \\
\end{longtable}}
